\subsection{All cuspidal abelian fibres in the full metabelian tripod}

\begin{statement}{A.172}{Proved}
{Exact cuspidal fibres and a uniform finite separator of the three labels}
Let \(\mathcal S^{\rm met}(U)\) be the geometric conjugacy classes
of continuous sections of \(E_U/\Pi''\) locally matching actual
point or cusp sections.  Write their full abelian coordinates as
\((A,B)\in(\widehat{\Q^\times})^2\), and put
\[
 \begin{split}
 \Omega(A,B)={}&\{(0,A):B=1\}\\
       &{}\cup\{(1,B):A=1\}
        \cup\{(\infty,A^{-1}):B=\kappa(-1)A\}.
 \end{split}
\tag{A.172.1}
\]
The coordinates for the three cusp parameters are respectively
\(t\), \(1-t\), and \(1/t\).
Whenever \(\Omega(A,B)\) is nonempty, the full fibre of
\(\mathcal S^{\rm met}(U)\to(\widehat{\Q^\times})^2\)
over \((A,B)\) is in bijection with \(\Omega(A,B)\), by the
images of the actual full cusp sections with these data.

Each such fibre has one or two elements.  The only overlaps of
the three cuspidal abelian loci are
\[
 \begin{array}{c|c}
 (A,B)&\Omega(A,B)\\\hline
 (1,1)&\{(0,1),(1,1)\}\\
 (\kappa(-1),1)&\{(0,\kappa(-1)),(\infty,\kappa(-1))\}\\
 (1,\kappa(-1))&\{(1,\kappa(-1)),(\infty,1)\}.
 \end{array}
\tag{A.172.2}
\]
These two-element fibres remain distinct in an explicit finite
geometric quotient \(Q^{(2)}\) of order \(27\) and nilpotency class
two.  It is obtained from the actual proper Fermat-cubic mod-three
quotient \(Q\), which has order \(81\) and class three.
More strongly, over \(\Q_5\) the quotient \(Q^{(2)}\)
separates any two actual cusp sections with different labels,
independently of their tangent parameters.  It therefore separates
the three global labels uniformly.  Class two is necessary for
this uniform distinction, since the first row of (A.172.2) already
gives two actual full cusp sections with equal full abelian image.

The full geometric pro-three quotient of class two separates the
three cusp labels after base change to every finite extension of
every \(\Q_p\), again for arbitrary tangent parameters.
Furthermore the entire cuspidal-fibre classification (A.172.1)--(A.172.2)
holds over every number field \(k\), with completed \(k^\times\),
the actual arithmetic extension over \(k\), and local matching
at its places.  The finite order-twenty-seven assertion above
retains its stated \(\Q_5\) scope.
\end{statement}

\begin{proof}
\emph{The candidate classes and their exact abelian maps.}
The rational local parameters give full tangential splittings of
the cusp decomposition extensions.  Twisting one by a continuous
Kummer cocycle gives an actual full cusp section for every element
of \(\widehat{\Q^\times}\).  The pullbacks of the two
global unit functions are exactly
\[
 \begin{array}{c|c|c}
 \text{cusp}&\text{parameter}&(A,B)\\\hline
 0&t&(\tau,1)\\
 1&1-t&(1,\tau)\\
 \infty&z=1/t&(\tau^{-1},\kappa(-1)\tau^{-1}).
 \end{array}
\tag{A.172.3}
\]
For example \(1-t=-z^{-1}(1-z)\) at infinity.  Each displayed
formal factor with constant term one has its unique formal
\(n\)-th root with constant term one in \(\Q[[z]]\), for
every \(n\); it contributes zero to all the finite Kummer
classes.  Thus (A.172.3) retains the sign and the actual maps,
not only their valuation vectors.  Solving these three equations
gives exactly (A.172.1)--(A.172.2).  There is no triple
intersection, since \(\kappa(-1)\ne1\), already modulo two
over \(\Q_2\).  For a fixed label its parameter is recovered
from (A.172.3), so different parameters have different full
abelian images.

\emph{The actual finite quotient of geometric order eighty-one.}
Choose geometric meridians \(x,y\) around zero and one, and set
\[
 H_3=\ker\bigl(\Pi\longrightarrow(\Z/3)^2\bigr).
\tag{A.172.4}
\]
Its arithmetic cover in the trivial Kummer class is
\[
       Y_3:\ U^3+V^3=1,\quad UV\ne0,
                 \qquad t=U^3,
       \qquad C_3:\ U^3+V^3=Z^3.
\tag{A.172.5}
\]
The first map is connected finite etale of degree nine.  The
proper curve is a smooth plane cubic and has genus one.  There
are three points above each of the three removed points, with
ramification index three.
Let \(K\subset H_3\) be the kernel of the actual
compactification and mod-three homology map
\[
         H_3\longrightarrow H_1(C_{3,\overline\Q},\F_3).
\tag{A.172.6}
\]
Equivalently \(K\) is the closed subgroup generated by
\([H_3,H_3]\), all cubes in \(H_3\), and the nine boundary
inertia groups.  This equality follows by abelianizing the
punctured genus-one presentation: the nine meridians have their
one sum relation, and quotienting them leaves exactly the two
compact homology generators.  Thus \(H_3/K\simeq(\Z/3)^2\)
and \([\Pi:K]=81\).

The subgroup \(K\) is normal in \(E_U\).  Indeed \(H_3\)
is characteristic in the geometric group, and each arithmetic
conjugation induces a semilinear map of its geometric cover.
That map extends to the proper normalization and permutes its
boundary.  It therefore preserves the exact kernel (A.172.6).
Consequently \(E_U/K\) is an actual arithmetic quotient,
with geometric group \(Q=\Pi/K\).  Moreover
\(\Pi'\subset H_3\) and
\(\Pi''\subset\overline{[H_3,H_3]}\subset K\), so this
quotient receives the full metabelian one.

Here is a literal presentation of \(Q\).  Write
\[
 \begin{gathered}
 R=\Z[\zeta_3]/(3(\zeta_3-1)),\qquad \delta=\zeta_3-1,\\
 R=(\Z/9)\oplus(\Z/3)\delta,
             \qquad 3\delta=0,\quad\delta^2=-3,\quad\delta^3=0,\\
 Q_0=R\rtimes\langle\zeta_3\rangle,
             \qquad (r,i)(s,j)=(r+\zeta_3^is,i+j),\\
 X=(0,1),\qquad Y=(1,1).
 \end{gathered}
\tag{A.172.7}
\]
The integer matrices of the relations \(3\delta=0\) and
\(3\delta^2=-9-9\delta=0\) have index twenty-seven, which
gives the stated additive group of \(R\).
Here \(X\) and \(Y\) generate: \(YX^{-1}\) is translation
by one, and conjugating it by \(X\) gives translation by
\(\zeta_3\).  The identity
\(1+\zeta_3+\zeta_3^2=0\) gives
\[
               X^3=Y^3=(XY)^3=1.
\tag{A.172.8}
\]
Its abelianization is \((R/\delta R)\oplus\Z/3\).
In the coordinates \((a+b\delta,i)\), the usual meridian
exponents are precisely
\[
              (a+b\delta,i)\longmapsto(i-a,a)\pmod3.
\tag{A.172.9}
\]
Thus the homomorphism from the free geometric group taking
\(x\) to \(X\) and \(y\) to \(Y\) maps \(H_3\) onto
\(\delta R\).  This subgroup is elementary abelian of order
nine.  Equations (A.172.8) kill every compactification meridian,
which is a conjugate of \(x^3\), \(y^3\), or \((xy)^3\).
Hence the map kills the exact kernel (A.172.6), and factors through
\(\Pi/K\).  Both this group and \(Q_0\) have order eighty-one,
so it is the claimed isomorphism, with both meridians specified.

Using commutators \([g,h]=ghg^{-1}h^{-1}\), the lower central
series is
\[
 \begin{gathered}
 [X,Y]=\delta,\qquad [[X,Y],X]=3\ne0,\\
 \Gamma_2Q=\delta R,\qquad
 \Gamma_3Q=\delta^2R,\qquad\Gamma_4Q=1,\\
                  (|Q|,|\Gamma_2Q|,|\Gamma_3Q|,|\Gamma_4Q|)
                             =(81,9,3,1).
 \end{gathered}
\tag{A.172.10}
\]
In particular this is a class-three quotient, with a literal
nonzero commutator of weight three.

\emph{Uniform separation of the cusp labels at the finite place five.}
On the proper curve (A.172.5), the actual elliptic isomorphism,
with origin \(P_\infty=(1:-1:0)\), is
\[
 \begin{gathered}
 E_0:\ Y_e^2=X_e^3-432,\qquad
 X_e=\frac{12Z}{U+V},\quad Y_e=\frac{36(U-V)}{U+V},\\
 \frac UZ=\frac{36+Y_e}{6X_e},\qquad
 \frac VZ=\frac{36-Y_e}{6X_e},\\
 P_0=(0:1:1)\mapsto(12,-36),\qquad
 P_1=(1:0:1)\mapsto(12,36),\qquad P_\infty\mapsto O.
 \end{gathered}
\tag{A.172.11}
\]
Substitution verifies both function-field maps.  They extend to
inverse maps of the proper smooth curves.  The tangent slope at
\((12,36)\) is six, so doubling gives \((12,-36)\); this
point has exact order three.  The three boundary images form its
three distinct multiples.

At five this model is smooth.  Its six residue points are
\[
      E_0(\F_5)=\{O,(1,2),(1,3),(2,1),(2,4),(3,0)\}.
\tag{A.172.12}
\]
The image of multiplication by three is \(\{O,(3,0)\}\).
Every nonzero difference of the three displayed boundary points
reduces to \((2,1)\) or \((2,4)\).  It is therefore not a
triple in \(E_0(\Q_5)\), since reduction would make it a
triple in (A.172.12).  The finite Kummer injection consequently
gives distinct classes
\[
             \kappa_3(P_0),\quad\kappa_3(P_1),\quad
                  \kappa_3(P_\infty)
                   \quad\text{in }H^1(\Q_5,E_0[3]).
\tag{A.172.13}
\]

Suppose two actual cusp sections at different labels had conjugate
images in \(E_{U,\Q_5}/K\).  Their abelian mod-three classes
would agree.  Since \(-1\) is a cube, the three mod-three cusp
loci in \(H^1(\Q_5,\mu_3)^2\) are
\[
                  (a,0),\qquad(0,b),\qquad(c,c).
\tag{A.172.14}
\]
Any two intersect only at zero.  Each finite section can therefore
be conjugated to the trivial Kummer section of (A.172.5).
The residual geometric deck centralizer is
\(\mu_3(\Q_5)^2=0\), since five is two modulo three.
Hence any further geometric conjugator between the two normalized
sections lies in \(H_3/K\), rather than merely in \(Q\).

The fixed-sheet lifts of the three local cusp labels have the
unique rational boundary centres \(P_0,P_1,P_\infty\) above
those labels.  These are literal formal lifts of
\(U^3=t\), \(V^3=1-t\): at infinity their ratio tends to
\(-1\), the unique cube root of minus one in \(\Q_5\).
Their projections to the proper quotient (A.172.6) are exactly
the three classes in (A.172.13).  Conjugation by \(H_3/K\)
changes their cocycles by a coboundary, contradicting that
inequality.  This proves the asserted local separation for all
tangent parameters, and by localization proves global separation.
In particular the two candidates in each row of (A.172.2) have
different images in this single quotient.

\emph{Exhausting every cuspidal abelian fibre.}
Fix \((A,B)\) with \(\Omega(A,B)\ne\varnothing\), and
let \(s\) be a locally matching full metabelian section with
that pair.  At every finite place it is cuspidal.  Indeed one of
\(B=1\), \(A=1\), or \(B=\kappa(-1)A\) holds globally,
and an ordinary local centre \(t\) would give respectively
\(t=0\), \(t=1\), or \(1-t=-t\), by full local Kummer
injectivity.  No one of these is an ordinary point of \(U\).

Choose one actual full cusp reference \(\sigma_0\) from
(A.172.1).  Lift an abelian conjugator and conjugate \(s\), so
its full abelian cocycle is literally the one of \(\sigma_0\).
Normalize the at most two full reference sections indexed by
\(\Omega(A,B)\) in this same way.  For \(n=6j!\) use the
actual subgroup tower and its arithmetic covers
\[
 \begin{gathered}
 H_n=\ker(\Pi\longrightarrow\Pi^{\rm ab}/n\Pi^{\rm ab}),
                \qquad H_n'=\overline{[H_n,H_n]},\\
 Y_n:\ A_nU^n+B_nV^n=1,\quad UV\ne0,
                    \qquad t=A_nU^n.
 \end{gathered}
\tag{A.172.15}
\]
Here \(A_n,B_n\in\Q^\times\) represent the two completed
classes modulo \(n\).  The arithmetic subgroup is the preimage
of the common finite section in \(E_U/H_n\).  The compatible
fixed sheets of \(\sigma_0\) specify these covers and their
transition maps.  For \(n\mid m\), those maps are
\[
  (U_m,V_m)\longmapsto(rU_m^{m/n},sV_m^{m/n}),\qquad
         r^n=A_m/A_n,\quad s^n=B_m/B_n.
\tag{A.172.16}
\]
They are the actual maps of the nested arithmetic subgroups.  In
particular their constants are rational: the indicated root-function
ratios have constant \(n\)-th powers, and the geometrically
connected curves have constant field \(\Q\).  This is the
same fixed-sheet construction as A.169.12--13.

The proper plane model \(C_n\) is smooth of genus
\((n-1)(n-2)/2\).  It has a rational boundary origin supplied
by \(\sigma_0\), with the explicit presentations
\[
 \begin{array}{c|c|c}
 \text{reference label}&\text{coefficient choice}&\text{rational boundary point}\\\hline
 0&B_n=1&(0:1:1)\\
 1&A_n=1&(1:0:1)\\
 \infty&B_n=-A_n&(1:1:0).
 \end{array}
\tag{A.172.17}
\]
All actual local lifts of \(s\) have proper centres in the
finite boundary \(Z_n=C_n\setminus Y_n\).  Because
\(\Pi''\subset H_n'\), the section gives an actual proper
inverse-Selmer class on \(\operatorname{Jac}(C_n)\).
Stoll's Theorem 3.11 applies to the Abel--Jacobi image of the
whole finite scheme \(Z_n\), with its nonrational points
retained \cite[Theorem 3.11]{LegStoll2007}.  It gives one
\(Q_n\in Z_n(\Q)\) with that class.  The Abel--Jacobi map
is injective on geometric points, and full local Kummer
injectivity makes every finite local centre equal to \(Q_n\).

Compare the affine abelian section on \(Y_n\) with a rational
tangential splitting at \(Q_n\).  Its difference is a class
in the actual boundary module
\(\widehat\Z(1)[Z_n]/\widehat\Z(1)\sum e_Q\), locally
in the line of \(Q_n\).  Compact Tate invariants vanish at
every finite place, so this difference is well-defined and its
boundary-module cohomology injection is exact.  The nonsplit
inertia-line recovery A.169.9--11 gives one global tangent
parameter.  It produces a full cusp section \(\sigma_n\)
on \(Y_n\) equal to the given section in its affine abelian
quotient.

Push \(\sigma_n\) through the actual finite cover to \(U\).
The formal map of local parameters is
\(z_i=c z_{Q_n}^{e}u(z_{Q_n})\), with \(c\in\Q^\times\)
and \(u(0)=1\); hence its full parameter is
\(c\tau_n^e\).  Its abelian pair is exactly \((A,B)\),
so (A.172.3) identifies this pushed full cusp class with one of
the fixed finite set \(\Sigma\) of references indexed by
\(\Omega(A,B)\).  No consistency of that choice as \(n\)
varies is presumed.

Instead retain both the reference and conjugator in one compact
space:
\[
 \mathcal C_n=\{(g,\sigma)\in(\Pi/\Pi'')\times\Sigma:
       g s(\gamma)g^{-1}=\sigma(\gamma)
       \text{ in }E_U/H_n'\text{ for every }\gamma\in G_\Q\}.
\tag{A.172.18}
\]
These sets are nonempty, closed, and nested under the actual
divisibility transitions.  Their intersection is nonempty because
\(\Sigma\) is finite and \(\Pi/\Pi''\) is compact.
The exact finite-character argument A.169.18 gives
\(\bigcap_jH_{6j!}'=\Pi''\).  Thus a single pair in the
intersection conjugates \(s\) to a single full reference cusp
in the full metabelian quotient.  Conversely every listed reference
is an actual full cusp section and is locally matching.  Same-label
parameters are separated by (A.172.3), and different labels by
(A.172.4)--(A.172.14).  This proves the claimed bijection.

\emph{The order-twenty-seven quotient and its exact central character.}
The detector above has a further arithmetic quotient
\[
 Q^{(2)}=Q/\Gamma_3Q
     \simeq\bigl(\F_3[\delta]/\delta^2\bigr)\rtimes C_3,
          \qquad \zeta_3\ \text{acts by }1+\delta.
\tag{A.172.19}
\]
The preimage of \(\Gamma_3Q\) in \(E_U\)'s geometric subgroup
is arithmetic normal because \(\Gamma_3Q\) is characteristic
in \(Q\).  Thus (A.172.19) specifies an actual quotient of
\(E_U/K\), of geometric order twenty-seven.  Its commutator is
the nonzero central line \(\delta\F_3\), so its class is two.

The meridian action on the first geometric abelian quotient is
multiplication by the cyclotomic character.  Bilinearity of
commutators on successive lower-central quotients therefore gives
the literal action \(\chi_3^3=\chi_3\) on
\(\Gamma_3Q=\langle[[X,Y],X]\rangle\).  Inner changes of the
arithmetic lift act trivially on this central subgroup.  Under the
proper comparison \(H_3/K=E_0[3]\), it is consequently the
nontrivial cyclotomic line in
\[
 E_0[3]=\langle P\rangle\oplus\langle T\rangle,\qquad
        P=(12,36),\quad T=(0,12\sqrt{-3}),\qquad
                 \Gamma_3Q=\langle T\rangle.
\tag{A.172.20}
\]
Both displayed points have order three by the tangent formula, and
they are independent since their nonzero multiples have different
\(X_e\)-coordinates.  The first line is rational and the second
has action \(\chi_3\); these characters are distinct, giving the
claimed identification.

The following actual isogenies give the projected Kummer class:
\[
 \begin{gathered}
 E_1:\ Y^2=X^3+11664,\\
 \varphi:E_0\longrightarrow E_1,\qquad
 (x,y)\longmapsto
   \left(\frac{x^3-1728}{x^2},
               \frac{y(x^3+3456)}{x^3}\right),\\
 \psi:E_1\longrightarrow E_0,\qquad
 (X,Y)\longmapsto
   \left(\frac{X^3+46656}{9X^2},
               \frac{Y(X^3-93312)}{27X^3}\right).
 \end{gathered}
\tag{A.172.21}
\]
Substitution verifies both curve equations and
\(\psi\varphi=[3]\).  The formulas extend over the proper
curves and send origins to origins.  Each has degree three, as
seen from the six poles of its target \(x\)-coordinate counted
with multiplicity.  Their kernels are respectively
\(\langle T\rangle\) and
\(\langle T_1\rangle\), \(T_1=(0,108)\), and
\(\varphi(P)=T_1\).
Thus \(\varphi\) identifies \(E_0[3]/\Gamma_3Q\) with the
constant group \(E_1[\psi]=\langle T_1\rangle\).
On cocycles the image of \(\kappa_3(S)\), for
\(S\in E_0(k)\), is exactly the \(\psi\)-connecting class
\(\partial_\psi(S)\): choose \(Z\) with \(3Z=S\), and
\(\varphi Z\) is a \(\psi\)-preimage whose coboundary is the
image under \(\varphi\) of the original Kummer cocycle.

For \(S=P\) a literal preimage is
\[
 \begin{gathered}
 Q_u=\bigl(36u,108(2u^2-2u-1)\bigr),\qquad
          u^3-3u^2+1=0,\qquad \psi(Q_u)=P,\\
 u=1-(\zeta_9+\zeta_9^{-1}),\qquad
 \sigma_2(u)=-u^2+2u+2,\qquad
             \sigma_2(Q_u)=Q_u+T_1 ,
 \end{gathered}
\tag{A.172.22}
\]
where \(\sigma_2(\zeta_9)=\zeta_9^2\).  The cubic is
irreducible by the rational-root test and has discriminant
eighty-one.  The cyclotomic expression and addition identity are
verified directly by substitution.  Consequently, with \(T_1\)
identified with \(1\in\F_3\), its exact connecting character is
\[
       \partial_\psi(P)(\gamma)
           =\vartheta_9(\gamma)
           :=\frac{\chi_9(\gamma)^2-1}{3}\pmod3.
\tag{A.172.23}
\]
The numerator is divisible by three, and changing its lift modulo
nine does not change the result.  The product rule modulo nine
proves this is an additive character.  It takes value one on
\(\sigma_2\), as in (A.172.22), and cuts out precisely the
real cyclic cubic field \(\Q(\zeta_9)^+\).

The three proper boundary classes project in \(Q^{(2)}\) to
\(-\vartheta_9,\vartheta_9,0\), respectively, under the origin
and normalization in (A.172.11).  At five,
\(\vartheta_9(\operatorname{Frob}_5)=2\ne0\).
They remain pairwise different there.  The same zero deck-centralizer
and fixed-sheet argument used for (A.172.13)--(A.172.14) applies
after this further quotient, whose abelian deck group is still
\((\Z/3)^2\).  Therefore the uniform separation already holds
in the actual order-twenty-seven quotient.  The two equal
full-abelian classes in the first row of (A.172.2) prove the
asserted necessity of class two.

\emph{The full pro-three class-two separator over every finite local
extension.}
Let \(\Pi_3\) be the maximal geometric pro-three quotient of
\(\Pi\), and take its actual class-two quotient
\[
 \begin{gathered}
 N_3=\Pi_3/\overline{\Gamma_3\Pi_3},\qquad
 L_3=\ker(\Pi\longrightarrow N_3),\qquad
 \mathcal E_{U,3}^{[2]}=E_U/L_3,\\
 1\longrightarrow M_3=\Z_3(2)\longrightarrow N_3
       \longrightarrow T_3=\Z_3(1)^2\longrightarrow1 .
 \end{gathered}
\tag{A.172.24}
\]
The last group is the free class-two pro-three group on the two
meridians: every element has the unique form
\(x^a y^b[x,y]^c\), \(a,b,c\in\Z_3\).  This gives the
displayed central sequence and its actual inertia maps.
Explicitly this form maps to the upper triangular unipotent
matrix with upper entries \(a,b,ab+c\); multiplication sends
\((a,b,c),(a',b',c')\) to
\((a+a',b+b',c+c'-a'b)\).
Collection of words and reduction modulo every \(3^n\)
give the inverse map, so this is the full pro-three group.
The arithmetic action on the centre is \(\chi^2\), by the
commutator calculation.  Sending the two meridians to those in
(A.172.19) gives the actual quotient \(N_3\to Q^{(2)}\);
its map on the centre is reduction modulo three, with the
nonzero generator specified by their commutator.

Take the three full tangent-one cusp sections over \(\Q\).
Their full pro-three abelian classes all agree, since the
pro-three Kummer class of \(-1\) is zero.  Lift their abelian
conjugators and normalize their abelian cocycles to that of the
infinity section \(\sigma_\infty\).  Then the literal ratios
\(\sigma_i(g)\sigma_\infty(g)^{-1}\) are continuous cocycles
with values in \(M_3\).  Since \(T_3^{G_\Q}=0\), the
normalizing abelian conjugator is unique; two lifts differ
centrally and change these ratios only by a central coboundary.
The exact calculation
A.40.1--2, whose domain is continuous cohomology of the full
absolute group, gives
\[
 H^1_{\rm cts}(G_\Q,M_3)=\Z/3,\qquad
              \eta_3(g)=\frac{\chi_{3^\infty}(g)^2-1}{24}.
\tag{A.172.25}
\]
Their three central classes are all different, because their
images in the actual quotient \(Q^{(2)}\) were separated above.
They are therefore precisely the three elements of this group;
every nonzero pairwise difference is \(\eta_3\) or \(2\eta_3\).
This use of A.40 concerns the difference of existing sections:
no lifting obstruction is being omitted.

Let \(K/\Q_p\) be any finite extension and use its actual
decomposition map \(G_K\to G_\Q\).  For \(k=1,2\), the
restriction of \(k\eta_3\) cannot be a coboundary in
\(H^1_{\rm cts}(G_K,\Z_3(2))\).  Such a coboundary would be
\((\chi_{3^\infty}^2-1)z\) for an integral \(z\in\Z_3\),
and hence would give
\[
 (\chi_{3^\infty}(g)^2-1)
                 \left(\frac{k}{24}-z\right)=0
                        \quad(g\in G_K).
\tag{A.172.26}
\]
The squared cyclotomic character is nontrivial on \(G_K\).
For \(p\ne3\), its value on arithmetic Frobenius is
\(|\kappa_K|^2>1\), as an ordinary rational integer in \(\Q_3\).
For \(p=3\), finite image would put all \(3\)-power roots of
unity in one finite extension of \(K\), which has only finitely
many roots of unity.  Choosing a nonzero coefficient in
(A.172.26) forces \(z=k/24\), whose three-adic valuation is
minus one.  This is impossible.  Thus restriction is injective
on the global group (A.172.25).

The same nontrivial cyclotomic action gives \(T_3^{G_K}=0\).
Consequently any geometric conjugator between two of the
normalized local sections in \(\mathcal E_{U,3}^{[2]}\)
has zero abelian projection and lies in \(M_3\).  It changes
their difference only by a central coboundary, contrary to
(A.172.26).  They are locally nonconjugate.

Finally consider arbitrary actual cusp parameters over \(K\).
Their abelian pro-three classes lie in the three coordinate
lines \((a,0),(0,b),(c,c)\), with
\(a,b,c\in H^1_{\rm cts}(G_K,\Z_3(1))\).
Two different labels can have conjugate images only when these
classes are all zero.  The corresponding cusp inertia cocycle
is then a coboundary, so conjugation inside its actual pro-three
inertia takes the image to the tangent-one section of that label.
The nonconjugacy just proved excludes this.  This establishes
the full local-field assertion, including the coefficient place
three and arbitrary finite ramification.

\emph{The exact number-field transfer.}
Fix a number field \(k\subset\overline\Q\).
With the same geometric base point, the actual arithmetic group
is the preimage of \(G_k\) in \(E_U\), namely
\(E_{U_k}=E_U\times_{G_\Q}G_k\).  Its geometric groups and all
displayed quotient maps are those already specified.  Interpret
(A.172.1) with \(A,B\in\widehat{k^\times}\) and the actual
Kummer class of minus one over \(k\).  This class is nonzero:
otherwise \(-1\) would have a \(2^m\)-th root in \(k\) for
every \(m\), giving roots of unity of unbounded order in that
number field.  The same three intersections are therefore
exactly (A.172.2).

Every step exhausting these fibres has the required number-field
map.  The reference cusp gives coefficients \(A_n,B_n\in k^\times\)
in (A.172.15), and the arithmetic subgroup is the preimage of its
common section in \(E_{U_k}/H_n\).  Its transition constants
in (A.172.16) lie in \(k\): their powers are constant, and the
geometrically connected function fields have constant field \(k\).
The three points in (A.172.17) are \(k\)-rational boundary
origins.  Full local Kummer injectivity over each finite
completion of \(k\) excludes ordinary centres exactly as above.
The resulting proper inverse-Selmer classes have their local
centres in the actual finite boundary scheme.
Stoll's Theorem 3.11 is stated over an arbitrary number field
and supplies one \(k\)-rational boundary point on each cover
\cite[Theorem 3.11]{LegStoll2007}.

For the residual boundary class, take a finite Galois splitting
field over \(k\).  The completed Kummer injection A.169.4 has
that number-field domain, and A.169.8--11 supplies both restriction
injectivity and descent of the single inertia parameter.
Thus the pushforward is one of the same finite set of full
\(k\)-cusp references.  The compact sets (A.172.18) are formed
in the unchanged geometric group and the finite candidate set,
with \(G_k\) in their equality conditions.  Their intersection
proves the same surjectivity onto the fibre.
For injectivity at different labels, restrict to any finite
completion of \(k\) and apply (A.172.24)--(A.172.26).
Same-label parameters remain separated by their full abelian
coordinates.  This proves the precise number-field corollary
without an additional Tate--Shafarevich hypothesis.

The ring, finite group, exact lower central series and meridian
maps, the two inverse cubic maps, all six nonzero local boundary
differences at five, the literal isogenies and cyclotomic preimage,
and the rational boundary presentations are
checked in \path{code/tripod_cuspidal_metabelian_certificate.py},
with retained output and exact data.  The full fibre assertion
uses the continuous cohomology and compactness argument above.
\end{proof}
